\documentclass{hoja-ejercicio}

% --- Macros propias del practico (vienen del .tex original) ---
\newcommand{\un}[1]{\,\mathrm{#1}}
\newcommand{\ms}{\un{m/s}}
\newcommand{\mss}{\un{m/s^2}}
\newcommand{\vv}[1]{\vec{#1}}
\newcommand{\ux}{\hat{\imath}}
\newcommand{\uy}{\hat{\jmath}}
\DeclareMathOperator{\sen}{sen}

\setpractico{Práctico 1 --- Cinemática 1D y 2D}
\setpiedepagina{Física 1 --- Lic. en Ciencias de la Computación --- 2026}

\begin{document}

\begin{hoja}[2]{15}

\begin{ejercicio}
En un partido de fútbol femenino, que se juega sobre piso de arena, una jugadora se encuentra frente al arco y nota que éste está desguarnecido. Patea la pelota con una velocidad inicial de módulo igual a $24\ms$, formando un ángulo de $30^\circ$ con la horizontal, intentando convertir el gol. Debido a que existe un fuerte viento, la pelota experimenta tanto la aceleración de la gravedad como una aceleración horizontal de $2\mss$ en sentido opuesto a su dirección de movimiento. Sabiendo que el arco tiene $2.44\un{m}$ de altura y está ubicado a $45\un{m}$ de donde parte la pelota, y que al impactar en la arena la pelota no se desplaza horizontalmente:
\begin{enumerate}[label=(\alph*), leftmargin=*, itemsep=1pt, topsep=3pt]
\item Dibuje un esquema de la situación y el sistema de coordenadas elegido.
\item Escriba los vectores aceleración, velocidad y posición de la pelota mientras la pelota está en vuelo.
\item Calcule la altura máxima que alcanza la pelota.
\item Determine si la jugadora logra convertir el gol. Realice todos los cálculos necesarios para justificar su respuesta.
\item Si no hubiese existido la aceleración horizontal provocada por el viento, ¿la jugadora habría convertido el gol? Realice todos los cálculos necesarios para justificar su respuesta.
\item En el esquema del ítem (a), grafique cualitativamente la trayectoria en el caso de que existe aceleración horizontal y dibuje los vectores velocidad y aceleración en dos puntos cualquiera de la misma.
\end{enumerate}
Considere que la aceleración de la gravedad es igual a $10\mss$.
\end{ejercicio}

\begin{solucion}

\apartado{(a) Esquema y sistema de coordenadas}
Origen en el punto donde se patea la pelota, sobre la arena; eje $x$ horizontal hacia el arco, eje $y$ vertical hacia arriba; $t=0$ en el puntapié. El arco es un segmento vertical en $x=45\un{m}$ entre $y=0$ y $y=2.44\un{m}$. El viento produce una aceleración constante hacia $-x$.

\begin{figuracol}
\begin{tikzpicture}
\begin{axis}[
    width=11.6cm, height=6.2cm,
    xmin=0, xmax=52, ymin=0, ymax=9,
    xtick={0,10,20,30,40,45,50}, ytick={0,2.44,5,7.2},
    xticklabel style={font=\tiny,/pgf/number format/fixed}, yticklabel style={font=\tiny,/pgf/number format/fixed},
    grid=both, grid style={dashed, gray!50},
    xlabel={$x$ [m]}, ylabel={$y$ [m]}, label style={font=\small},
    legend style={font=\scriptsize, at={(0.97,0.97)}, anchor=north east},
    axis lines=left,
]
\addplot[domain=0:2.4, samples=60, hojaAzul, line width=1.5pt] ({20.785*x - x^2},{12*x - 5*x^2});
\addlegendentry{con viento ($a_x=-2\mss$)}
\addplot[domain=0:2.4, samples=60, hojaRojo!80!black, dashed, line width=1.2pt] ({20.785*x},{12*x - 5*x^2});
\addlegendentry{sin viento}
\draw[line width=2.5pt, black] (axis cs:45,0) -- (axis cs:45,2.44);
\node[font=\scriptsize, above right] at (axis cs:45,2.6) {arco};
\draw[->, thick, hojaAzul] (axis cs:8.15,4.0) -- (axis cs:16.1,7.2) node[right, font=\scriptsize] {$\vv v$};
\draw[->, thick, hojaRojo!80!black] (axis cs:8.15,4.0) -- (axis cs:7.35,0) node[pos=0.5, left, font=\scriptsize] {$\vv a$};
\draw[->, thick, hojaAzul] (axis cs:23.5,7.2) -- (axis cs:30.9,7.2) node[right, font=\scriptsize] {$\vv v$};
\draw[->, thick, hojaRojo!80!black] (axis cs:23.5,7.2) -- (axis cs:22.7,3.2) node[below, font=\scriptsize] {$\vv a$};
\addplot[only marks, mark=*, mark size=1.6pt, black] coordinates {(8.15,4.0) (23.5,7.2) (44.12,0) (45,2.54)};
\end{axis}
\end{tikzpicture}
\end{figuracol}

\apartado{(b) Vectores}
Condiciones iniciales: $\vv r_0=\vv 0$ y
\[
v_{0x} = 24\cos 30^\circ = 12\sqrt3\approx 20.78\ms,
\]
\[
v_{0y} = 24\sen 30^\circ = 12\ms .
\]
La aceleración es constante en ambos ejes (gravedad hacia abajo y viento hacia $-x$), así que integro dos veces componente a componente:
\begin{align*}
\vv a(t) &= -2\,\ux - 10\,\uy,\\
\vv v(t) &= (12\sqrt3 - 2t)\,\ux + (12-10t)\,\uy,\\
\vv r(t) &= (12\sqrt3\,t - t^2)\,\ux + (12t - 5t^2)\,\uy .
\end{align*}

\apartado{(c) Altura máxima}
Ocurre cuando $v_y=0$: $12-10t=0\Rightarrow t=1.2\un{s}$.
\[
y_{\max} = 12(1.2) - 5(1.2)^2 = 7.2\un{m}.
\]
En ese instante $x=20.78\cdot1.2-1.44\approx 23.5\un{m}$.

\apartado{(d) ¿Es gol con viento?}
Primero veo cuánto dura el vuelo: $y=0 \Rightarrow t(12-5t)=0 \Rightarrow t_v = 2.4\un{s}$. La pelota toca la arena en
\[
x(2.4) = 12\sqrt3\cdot 2.4 - 2.4^2 \approx 44.1\un{m} < 45\un{m}.
\]
Cae $0.9\un{m}$ antes de la línea del arco y, según el enunciado, al impactar no sigue avanzando. Otra forma de verlo: para llegar a $x=45$ haría falta $t^2-12\sqrt3\,t+45=0\Rightarrow t\approx 2.46\un{s} > t_v$, instante en que $y\approx -0.7\un{m}$ (ya estaría bajo el suelo).

\resultado{\textbf{No hay gol}: la pelota se queda corta, cae a $44.1\un{m}$ del punto de partida.}

\apartado{(e) ¿Y sin viento?}
Ahora $a_x=0$ y el eje $x$ es un MRU: $x(t)=12\sqrt3\,t$. El tiempo de vuelo no cambia ($2.4\un{s}$, depende sólo del eje $y$). Llega a la línea del arco en
\[
t_{\text{arco}} = \frac{45}{12\sqrt3}\approx 2.165\un{s} < 2.4\un{s},
\]
o sea que sigue en el aire, y a esa altura
\[
y(2.165) = 12(2.165) - 5(2.165)^2 \approx 2.54\un{m} > 2.44\un{m}.
\]

\resultado{\textbf{Tampoco es gol}: pasa unos $10\un{cm}$ \emph{por encima} del travesaño. El viento acorta el alcance de $49.9\un{m}$ a $44.1\un{m}$; para convertir haría falta una situación intermedia.}

\apartado{(f) Trayectoria con viento y vectores}
La curva ya no es una parábola simétrica: como $v_x$ disminuye durante el vuelo, la subida es más ``tendida'' y la bajada más empinada; el punto más alto está más cerca del punto de caída que del de partida ($23.5\un{m}$ de $44.1$). En el gráfico de (a) están dibujados los vectores en dos puntos:
\begin{itemize}[leftmargin=*, itemsep=1pt, topsep=3pt]
\item En $t=0.4\un{s}$ (subiendo, en $x\approx 8.2\un{m}$, $y=4\un{m}$): $\vv v=(20.0,\,8)\ms$ apunta hacia arriba y adelante, tangente a la curva; $\vv a=(-2,-10)\mss$ hacia abajo y hacia atrás. Es el mismo $\vv a$ en todo el vuelo.
\item En el punto más alto ($t=1.2\un{s}$): $\vv v = 18.4\,\ux\ \ms$, horizontal; $\vv a=(-2,-10)\mss$. A diferencia del tiro parabólico común, acá $\vv v$ y $\vv a$ no son perpendiculares (forman $\approx 101^\circ$) porque la componente $a_x$ sigue frenando la pelota.
\end{itemize}

\end{solucion}

\end{hoja}

\end{document}
